\subsection{The large-prime tail}\label{subsec:v5-large-sieve}

We use Lemmas~2.1 and 4.1 of Browning--Heath-Brown~\cite{v5bhb}.
The adjacent-norm polynomials are inhomogeneous, so we give the
fixed-coordinate argument rather than invoke their homogeneous
quadric-sieve theorem directly.
Their Lemma~2.1 says that coprime polynomials in $\mathbb Z[X_1,\ldots,X_m]$
of bounded degree and coefficient height at most $H$ have at most
\begin{equation}\label{eq:v5-box-sieve}
 O\left(\frac{V\log(VH)}{L\log L}
           +\frac{V\log(VH)}{B_{\min}}\right)
\end{equation}
integer points in a box of side parameters $B_i\geq2$ at which a prime
greater than $L\geq2$ divides every polynomial; $V=\prod B_i$.
``Coprime'' here excludes a common constant prime factor as well as a
common nonconstant factor. Their Lemma~4.1 covers the zeros modulo a
squarefree $q$ of a rank-$m$ quadratic form, with $m\geq3$, by at most
$(3m)^{\omega(q)}q^{m-2}$ lattices of determinant $q^{m-1}$.
The number with largest Euclidean successive minimum exceeding $T$ is
\begin{equation}\label{eq:v5-lattice-tail}
 O_m((3m)^{\omega(q)}q^{2m-3}T^{-m}).
\end{equation}
Every prime dividing $q$ must be odd and preserve the rank of the form.
The lattices in that lemma contain $q\mathbb Z^m$; in particular their
largest successive minima are at most $q$.

\begin{lemma}[Large common divisors on the quadric]\label{lem:v5-large-sieve}
Fix the implied constant in the box of size $O(B)$.
For fixed $\kappa\ne0$ and $\epsilon,\eta\in\{1,-1\}$, uniformly for
$B,L\geq2$, the number of integral vectors in this box satisfying
$Q_\kappa=0$ and
$\ell\mid G_\epsilon,H_\eta$ for some prime $\ell>L$ is
\begin{equation}\label{eq:v5-large-sieve}
 O_{\kappa,\alpha}\left(
    \frac{B^{3+\alpha}}{L\log L}+B^{11/4+\alpha}\right)
 \qquad(\alpha>0).
\end{equation}
\end{lemma}
\begin{proof}
Sign changes of $a,c$ reduce to $\epsilon=\eta=1$.
Put $x_0=a+c$, $x_1=a-c$, $x=b-d$, $z=b+d$, and
$R=\kappa w^2-xz$. Then $x_0x_1=R$.
Parity restrictions can be discarded in an upper bound.
For a fixed nonzero value of either $x_0$ or $x_1$, the triple $(x,z,w)$
determines the other coordinate uniquely whenever it gives an integral
solution. We shall use this injection when discarding divisibility
conditions.

First suppose that $\ell\mid x_0,x_1$. Then $\ell^2\mid R$.
For each fixed integer $v=O(B^2)$, the equation $R=v$ has
$O_{\kappa,\alpha}(B^{1+\alpha})$ triples: for each $w$, count the
factorizations $xz=\kappa w^2-v$ by the divisor bound; the zero right-hand
side occurs for at most two $w$ and contributes $O(B)$ each.
If $R\ne0$, it has $O_\alpha(B^\alpha)$ factorizations $x_0x_1=R$.
Since $0<|h|\ll B^2/\ell^2$ and $\ell\ll B$, summation gives
\[
 O_{\kappa,\alpha}\left(
 B^{3+2\alpha}\sum_{\ell>L}\ell^{-2}\right)
 =O_{\kappa,\alpha}(B^{3+2\alpha}/(L\log L)).
\]
The same bound for triples with $R=0$, followed by the $O(B)$ choices
with $x_0x_1=0$, gives $O_{\kappa,\alpha}(B^{2+\alpha})$.
We may discard all these $R=0$ solutions before proceeding.

In the remaining cases choose $t=x_1$ or $t=x_0$ such that $\ell\nmid2t$.
Take $t=x_1$ first. Here
$G_1-H_1=\kappa w^2+2t$ and $4t^2G_1=K$.
Thus the two norm congruences are equivalent to divisibility of
\[
 P=\kappa w^2+2t,\qquad
 K=(\kappa w^2-xz+t^2+2t)^2+t^2(x+z)^2.
\]
The identity
\[
 K-(t^2+x^2)(t^2+z^2)=P(P+2t^2-2xz)
\]
shows that it suffices to sieve the coprime pair
$P,(t^2+x^2)(t^2+z^2)$. They have no common nonconstant factor over
$\mathbb Q$, because the first depends nontrivially only on $w$, whereas
the second is independent of $w$. Their contents have no common prime,
since the second has coefficient $1$ on $x^2z^2$.
For $t=x_0$ we still have $4t^2G_1=K$, whereas
$t(G_1-H_1)=\kappa(t+2)w^2-2xz$. The pair is therefore
\[
 P'=\kappa(t+2)w^2-2xz,\qquad K.
\]
If $t\ne-2$, $P'$ is an irreducible ternary quadratic form of rank three
and $K$ has nonzero constant term $t^2(t+2)^2$; homogeneity rules out
divisibility. If $t=-2$, neither $x$ nor $z$ divides $K$.
The coefficient $1$ of $x^2z^2$ in $K$ excludes every common content
prime. In both charts the fixed value $t$ is nonzero.

In either chart, $t\mid R$. Put
\[
 S_R=\{p:p\mid2\kappa\},\qquad
 q=\prod_{\substack{p\mid t\ p\notin S_R}}p.
\]
The matrix of $R$ has determinant $-\kappa/4$, so $R$ has rank three
at every prime in this product, and those primes are odd. Also $q\ll B$.
For a fixed $q$, every eligible nonzero $t=O(B)$ is supported on
$S_R\cup\{p:p\mid q\}$. For every $\delta>0$, Rankin's bound gives
\[
 \#\{t:q(t)=q,\ 0<|t|\ll B\}
 \ll B^\delta
   \prod_{p\in S_R\ \text{or}\ p\mid q}(1-p^{-\delta})^{-1}
 \ll_{\kappa,\delta} B^\delta q^\delta
 \ll_{\kappa,\delta} B^{2\delta}.
\]
For the middle inequality, remove the finitely many primes with
$p^{-\delta}>1/2$ and bound the remaining factors by $2^{\omega(q)}
\ll_\delta q^\delta$. Renaming the small exponent gives
$O_{\kappa,\alpha}(B^\alpha)$ choices of $t$.

Apply \eqref{eq:v5-lattice-tail} with $m=3$ to $q\mid R$.
If $\lambda_1\leq\lambda_2\leq\lambda_3$ are the successive minima,
then $\lambda_3\leq q\ll B$ and
$\lambda_1\lambda_2\lambda_3\asymp q^2$.
A basis comparable to these minima therefore gives
$O(\prod_i(1+B/\lambda_i))=O(B^3/q^2)$ points in the box.
The lattices with $\lambda_3>T$ contribute at most
\[
 B^\alpha\sum_{q\ll B}
 O\left(9^{\omega(q)}q^3T^{-3}\frac{B^3}{q^2}\right)
 \ll_{\kappa,\alpha} B^{5+2\alpha}T^{-3}.
\]
For the other lattices choose a basis with lengths comparable to the
successive minima. The inverse basis matrix bounds the new coordinates by
$O(B/\lambda_i)$, so the resulting rectangular box has volume
$O(B^3/q^2)$ and minimum side parameter $\gg B/T$.
Enlarge its side parameters to at least $2$, as required in
\eqref{eq:v5-box-sieve}; since $\lambda_i\ll B$, this changes volume
only by a fixed factor.
The substitution into the polynomial pair is invertible over $\mathbb Q$.
For $\ell\nmid2t$ its determinant, of absolute value $q^2$, is invertible
modulo $\ell$ too. Thus no new common content prime outside $2t$ appears.
Remove common content supported on $2t$, which cannot affect the primes
being counted. Coefficient heights are powers of $B$ with fixed exponents.
Formula \eqref{eq:v5-box-sieve} bounds this lattice's contribution by
\[
 O\left(\frac{B^3\log B}{q^2L\log L}
                 +\frac{B^2T\log B}{q^2}\right).
\]
There are at most $9^{\omega(q)}q$ lattices. Since
$9^{\omega(q)}\ll_\alpha q^\alpha$, summing over $q\ll B$ and the
$O_{\kappa,\alpha}(B^\alpha)$ values of $t$ gives
$O_{\kappa,\alpha}(B^{3+3\alpha}/(L\log L)+B^{2+3\alpha}T)$.
Choose $T=B^{3/4}$ and rename the small exponent to obtain
\eqref{eq:v5-large-sieve}.
\end{proof}

\begin{proposition}[Weighted obstruction tail]\label{prop:v5-obstruction-tail}
For fixed distinct $c_j$ there are $L_*\geq2$ and $\tau_0,C>0$
such that, for every fixed $0<\tau\leq\tau_0$, uniformly for $L\geq L_*$,
the logarithmically weighted number of pairs
\[
 N/2\leq u\leq N,\qquad 1\leq w\leq\tau\sqrt N,
 \qquad p_j=u+c_jw^2\text{ all prime},
\]
having any of the odd common divisors forbidden in
Theorem~\ref{thm:v5-simultaneous-primes} greater than $L$, is at most
\begin{equation}\label{eq:v5-obstruction-tail}
 \frac{C\tau N^{3/2}}{L\log L}
       +O_\tau(N^{3/2-1/800}).
\end{equation}
The weight is $\prod_j\log p_j$, and the bound allows every Gaussian
generator, so it also holds for any prescribed choice of generators.
\end{proposition}
\begin{proof}
Take $B=\sqrt N$ and $L_0=B^{1/100}$, and choose $\tau_0$ so that
every prime value in the displayed region is between fixed positive
multiples of $B^2$. Thus $\prod_j\log p_j\ll (\log B)^k$.
Choose $L_*$ above the fixed exceptional primes in the medium-prime bound.
When $L\leq L_0$, that bound gives
\[
 \frac{C\tau B^3}{L\log L}+O_\tau(B^{27/10}(\log B)^k)
\]
for $L<\ell\leq L_0$. Its proof is uniform in this interval of $L$:
the remainder sum is bounded by the sum over all $\ell\leq L_0$.
When $L>L_0$ this part is empty.

For distinct positions, apply Lemma~\ref{lem:v5-large-sieve} at $L_0$
with $\alpha=1/1000$. After multiplication by the weight, the bound is
\[
 O\left(\left(
   \frac{B^{3+1/1000}}{B^{1/100}\log L_0}
       +B^{11/4+1/1000}\right)(\log B)^k\right)
 =O(B^{3-1/200}).
\]
This bounds all the remaining primes, also when $L>L_0$.
At one position, an odd divisor of both adjacent norms divides
$a$ and $b^2+1$. The coprime pair $a,b^2+1$ in variables $(a,b,w)$,
inserted in \eqref{eq:v5-box-sieve} at $L_0$, gives
\[
 O\left(\left(
 \frac{B^3\log B}{L_0\log L_0}+B^2\log B\right)(\log B)^k\right)
 =O(B^{3-1/200}).
\]
No factor $\tau$ is needed in these power-saving errors.

For a fixed pair of positions, the vector $(a,b,c,d,w)$ determines
$u=a^2+b^2-c_jw^2$; at one position $(a,b,w)$ does so.
Thus the vector counts bound the pattern counts. There are finitely many
pairs of positions and signs. Each split rational prime has eight
Gaussian elements of that norm, so counting all generator choices
changes the bound by at most a fixed factor. Every vector used in the
medium-prime count is primitive, since a common integer divisor greater
than one would have its square divide a prime norm.
All logarithmic factors in the remainders are absorbed for sufficiently
large $B$, with thresholds depending only on the fixed data and $\tau$.
Weakening the error to $B^{3-1/400}=N^{3/2-1/800}$ proves the claim.
\end{proof}
